% 融合算子多核 MatMul tiling 策略指南（AscendC / Atlas A2）
% 编译：xelatex -interaction=nonstopmode 融合算子多核tiling策略指南.tex（两遍）
\documentclass[UTF8,a4paper,11pt]{ctexart}
\usepackage{amsmath,amssymb}
\usepackage{booktabs}
\usepackage{geometry}
\usepackage{hyperref}
\usepackage{longtable}
\usepackage{listings}
\geometry{margin=2.2cm}
\hypersetup{colorlinks=true,linkcolor=blue,urlcolor=blue}
\lstset{basicstyle=\ttfamily\small,breaklines=true,frame=single}

\title{融合算子多核 Tiling 策略指南\\{\large AscendC KernelLaunch · Atlas A2 · 以 int8 MatMul 实验为据}}
\author{ascendc 工程}
\date{2026-06-11}

\begin{document}
\maketitle

\begin{abstract}
本文档从 \texttt{ascendc-tests/int8-matmul-cube-128x512x512/} 的多核 CPU 实验归纳
\textbf{Host 侧 MatMul tiling} 与 \textbf{Launch} 的闭合条件，并扩展到
\textbf{Cube + Vector 融合算子}（如 MatmulLeakyRelu）的配参思路。
实验平台为 CANN 9.0.0、Atlas A2（910B4）；讨论纪要见
\texttt{qa/2026-06/2026-06-11-ascendc-engineering-notes与数据搬运.md}。
\end{abstract}

\tableofcontents
\newpage

\section{适用范围与术语}

\subsection{适用范围}
\begin{itemize}
  \item \textbf{范式}：KernelLaunch；算子在 \texttt{*\_custom.cpp}，\texttt{main.cpp} 为应用壳。
  \item \textbf{硬件}：Atlas A2（AIC:AIV = 1:2）；\textbf{不}覆盖 Ascend 310P 等特例。
  \item \textbf{矩阵乘段}：\texttt{MultiCoreMatmulTiling} + 高阶 \texttt{Matmul<>}（纯 Cube 或融合 Vector 段）。
  \item \textbf{dtype 示例}：int8 $\times$ int8 $\rightarrow$ int32（实验）；fp16 等需替换最小 tile 与 base 倍数。
\end{itemize}

\subsection{两层分块}
\begin{table}[h]
  \centering
  \begin{tabular}{lll}
    \toprule
    层级 & Host 接口 & 决定内容 \\
    \midrule
    宏分片 & \texttt{SetDim}、\texttt{SetSingleShape}、\texttt{GetTiling} & 每个 \texttt{blockIdx} 负责 $C$ 的子块 \\
    微分片 & \texttt{SetFixSplit(baseM, baseN, baseK)} & 每个 AIC 内部 L0/Cube 迭代粒度 \\
  \end{tabular}
\end{table}

\textbf{常见误区}：认为只改 \texttt{blockDim} 与 \texttt{SetDim} 即可扩核。
宏分块数由 \texttt{GetTiling}（或 \texttt{SetSingleShape} 闭合公式）决定，与 launch 数必须一致。

\subsection{Launch 与 \texttt{usedCoreNum}（A2）}
\begin{itemize}
  \item \texttt{main.cpp} 中 \texttt{blockDim}：本次 launch 的 \textbf{AI Core（AIC）个数}。
  \item \texttt{SetDim(usedCoreNum)}：参与 MatMul 的 \textbf{AIV 规模}；本工程取
        $\texttt{usedCoreNum} = \texttt{blockDim} \times 2$。
  \item 内核用 \texttt{GetBlockIdx()}（0 $\ldots$ \texttt{blockDim}$-1$）配合
        \texttt{CalcGMOffset} 计算 GM 偏移。
\end{itemize}

\section{Tiling 闭合条件（用户定稿）}

以下在 \textbf{step/depth 不手改}、\textbf{baseK = -1}（K 向微块由 API 选择）、
\textbf{SingleCoreK = K}（宏上 K 不切，利于搬运与融合段衔接）时成立。

\subsection{硬件与 base}
\begin{enumerate}
  \item 实验设备 int8 入 / int32 出时，Cube 最小 tile 约为
        $16 \times 32 \times 16$（$M\times N\times K$）。
  \item 故 $\texttt{baseM}$ 为 \textbf{16 的倍数}，$\texttt{baseN}$ 为 \textbf{32 的倍数}。
  \item $\texttt{SetFixSplit(baseM, baseN, -1)}$；\texttt{baseK=-1} 不表示 K 向宏切分。
\end{enumerate}

\subsection{宏块与微块}
\begin{enumerate}
  \setcounter{enumi}{3}
  \item $\texttt{SingleCoreM} \bmod \texttt{baseM} = 0$，
        $\texttt{SingleCoreN} \bmod \texttt{baseN} = 0$。
  \item L0A/L0B 须能容纳当前 base 与 SingleCore 下的搬入策略（失败时常表现为
        \texttt{GetTiling} 返回 $-1$ 或 KFC 初始化错误）。
\end{enumerate}

\subsection{多核网格（核心等式）}
\begin{equation}
  \frac{M}{\texttt{SingleCoreM}} \times \frac{N}{\texttt{SingleCoreN}}
  = \texttt{usedCoreNum}
  \label{eq:grid}
\end{equation}
要求 $M,N$ 与 $\texttt{SingleCoreM,N}$ \textbf{整除}。
等价形式：
$\texttt{SingleCoreM} \times \texttt{SingleCoreN} \times \texttt{usedCoreNum} = M \times N$。
\textbf{$K$ 不参与式 \eqref{eq:grid}}。

\subsection{配参流程}
\begin{enumerate}
  \item 选定 $\texttt{usedCoreNum}$（= 宏分块总数 $= m_{\mathrm{blk}} \times n_{\mathrm{blk}}$）。
  \item 将 $\texttt{usedCoreNum}$ 因子分解，赋值
        $m_{\mathrm{blk}} = M/\texttt{SingleCoreM}$，
        $n_{\mathrm{blk}} = N/\texttt{SingleCoreN}$。
  \item 计算 $\texttt{SingleCoreM} = M/m_{\mathrm{blk}}$，
        $\texttt{SingleCoreN} = N/n_{\mathrm{blk}}$，$\texttt{SingleCoreK} = K$。
  \item 校验对 base 的倍数关系与 L0。
  \item Host：\texttt{SetSingleShape(SingleCoreM, SingleCoreN, K)}（推荐显式闭合），
        \texttt{SetDim(usedCoreNum)}，\texttt{SetFixSplit}。
  \item Launch：$\texttt{blockDim} = \texttt{usedCoreNum}/2$（A2）。
  \item \texttt{GetTiling} 后建议断言 $\texttt{res} \neq -1$，并核对
        \texttt{TCubeTiling} 中 \texttt{singleCoreM/N} 与预期一致。
\end{enumerate}

\section{实验用例：$128 \times 512 \times 512$ int8 MatMul}

\subsection{符号}
$A \in \mathbb{Z}^{128 \times 512}$，$B \in \mathbb{Z}^{512 \times 512}$，
$C \in \mathbb{Z}^{128 \times 512}$，int8 $\times$ int8 $\rightarrow$ int32。

\subsection{已验证配置摘要}
\begin{table}[h]
  \centering
  \small
  \begin{tabular}{ccccp{4.5cm}}
    \toprule
    AIC 数 & usedCoreNum & SetFixSplit & SetSingleShape & 结果 \\
    \midrule
    1 & 2 & (16,32,*) & 无 & 通过 \\
    4 & 8 & (16,32,-1) & 无 & 通过 \\
    8 & 16 & (32,32,-1) & 无 & 失败（GetTiling / SIGFPE） \\
    8 & 16 & (16,32,-1) & 无 & 通过（自动约 8 路 N 切） \\
    16 & 32 & (16,32,-1) & 无 & 失败（launch 16，tiling 约 8 块） \\
    16 & 32 & (16,32,-1) & (32,64,K) 等三种 & 通过 \\
    \bottomrule
  \end{tabular}
\end{table}

\subsection{16 AIC 下三种等价 SetSingleShape}
$\texttt{usedCoreNum}=32$，$\texttt{blockDim}=16$，均满足式 \eqref{eq:grid}：
\begin{itemize}
  \item $(32,64,K)$：$4 \times 8 = 32$
  \item $(64,32,K)$：$2 \times 16 = 32$
  \item $(16,128,K)$：$8 \times 4 = 32$
\end{itemize}
差异在 M/N \textbf{切分几何}与缓存局部性，非数学结果。

\subsection{反例：未闭合的 SetSingleShape}
$(128,32,K)$：$1 \times 16 = 16 \neq 32$ $\Rightarrow$ \texttt{blockIdx}$\geq 8$ 时
\texttt{tailN} 为负，\texttt{CheckTails} 断言失败。

\subsection{无 SetSingleShape 时的并行上界}
对 $N=512$、$\texttt{baseN}=32$，N 向理论最大块数
$\lfloor N / \texttt{baseN} \rfloor = 16$；但 \texttt{GetTiling} 在
本实验中 \textbf{自动} 常停在约 8 路（$\texttt{singleCoreN}=64$）。
欲稳定达到 16 路，需按 §2 显式 \texttt{SetSingleShape} 闭合。

\section{融合算子扩展（Matmul + Vector）}

\subsection{与纯 MatMul 的相同点}
\begin{itemize}
  \item MatMul 段仍遵守 §2 全部条件；\texttt{REGIST\_MATMUL\_OBJ}、
        \texttt{CalcGMOffset}、\texttt{SetTail}（或 device 侧 \texttt{SetSingleShape}）
        与 \texttt{TCubeTiling} 一致。
  \item A2 上 Cube:Vector = 1:2；\texttt{baseM}:\texttt{baseN} 宜为 1:2（如 16:32），
        与 int8 最小 tile 一致。
\end{itemize}

\subsection{融合段额外注意}
\begin{itemize}
  \item \textbf{VECIN/VECOUT} 队列深度与 \texttt{stepM/stepN}：实验约定 \textbf{不手改}
        除非整体验证；改 step 会牵动 Cube--Vector 流水线。
  \item Vector 段按 \texttt{singleCoreM} $\times$ \texttt{singleCoreN} 子块消费 Cube 输出；
        宏分片错误时 Cube 与 Vector 同步错位。
  \item 多 AIC 融合算子在 CPU 孪生上可能遇跨核 Matmul 同步问题；优先保证
        §2 闭合，再查融合 API；避免未闭合 tiling 时用多次 launch 掩盖根因（除非明确为 CPU 绕行）。
\end{itemize}

\subsection{推荐 Host 模板（伪代码）}
\begin{lstlisting}
// 1. 定 usedCoreNum, 因子分解, 算 SingleCoreM/N
tilingApi.SetDim(usedCoreNum);
tilingApi.SetOrgShape(M, N, K);
tilingApi.SetShape(M, N, K);
tilingApi.SetSingleShape(singleCoreM, singleCoreN, K);
tilingApi.SetFixSplit(baseM, baseN, -1);
tilingApi.SetBufferSpace(-1, -1, -1);
if (tilingApi.GetTiling(tilingData) == -1) { /* abort */ }
// 2. launch
blockDim = usedCoreNum / 2;  // A2
\end{lstlisting}

\section{故障诊断}

\begin{table}[h]
  \centering
  \small
  \begin{tabular}{lp{8cm}}
    \toprule
    现象 & 优先检查 \\
    \midrule
    \texttt{Ceiling} SIGFPE & \texttt{GetTiling} 是否 $-1$；\texttt{singleCoreM} 是否为 0 \\
    \texttt{tailN} 负 / \texttt{CheckTails} & 式 \eqref{eq:grid}；\texttt{blockDim} 是否大于宏块数 \\
    \texttt{totalIter\_}=0 & 该 \texttt{blockIdx} 无 N 向合法子块 \\
    部分 Block 成功 & launch 块数 $>$ \texttt{Ceiling(M,sM)*Ceiling(N,sN)} \\
    改 base 后高核数才过 & 微块约束导致 \texttt{GetTiling} 在高并行度下无解 \\
    \bottomrule
  \end{tabular}
\end{table}

\section{参考路径}

\begin{itemize}
  \item 实验目录：\texttt{ascendc-tests/int8-matmul-cube-128x512x512/}
  \item 讨论纪要：\texttt{qa/2026-06/2026-06-11-ascendc-engineering-notes与数据搬运.md}
  \item 模板：\texttt{MatmulInvocationNeo}、\texttt{MatmulLeakyReluInvocation}
  \item 平台纪要 §7.5：\texttt{qa/2026-06/2026-06-09-AscendC平台与CANN文档索引.md}
\end{itemize}

\end{document}
